\documentclass[11pt]{article}

\usepackage[margin=1in]{geometry}
\usepackage{amsmath,amssymb,amsthm}
\usepackage{booktabs}
\usepackage{algorithm}
\usepackage{algpseudocode}
\usepackage{natbib}
\usepackage{hyperref}
\usepackage{xcolor}
\usepackage{enumitem}

\hypersetup{colorlinks=true, linkcolor=blue, citecolor=blue, urlcolor=blue}

\newtheorem{theorem}{Theorem}
\newtheorem{proposition}{Proposition}
\newtheorem{definition}{Definition}
\newtheorem{assumption}{Assumption}

\newcommand{\neff}{n_{\mathrm{eff}}}
\newcommand{\todo}[1]{\textcolor{red}{\textbf{[TODO --- #1]}}}

\title{GT-Judge: A Game-Theoretic Multi-Agent Judge\\ for Bias-Resistant LLM Evaluation}

\author{Sawan Ade\\ IISER Bhopal \& IIT Patna\\ \texttt{sawan.ade@iiserb.ac.in}}

\date{July 18, 2026}

\begin{document}
\maketitle

\begin{abstract}
We introduce GT-Judge (Game-Theoretic Multi-Agent Judge), a training-free
evaluation framework that models the judging process as the outcome of a
three-stage strategic game pipeline. Traditional LLM-as-a-judge systems suffer
from human-like biases---position bias, verbosity bias, and sycophancy---that
are difficult to resolve via heuristic prompt tuning or majority voting.
GT-Judge targets these biases by coupling (1) an extensive-form Reasoning
Debate Game with endogenous claim validity and a token penalty to discourage
verbose or sycophantic arguments, (2) a Confidence Calibration Game using a
proper scoring rule (Brier reward) to elicit honest confidence estimates, and
(3) a Consensus Arbitration Game that resolves evaluator disagreements via the
Nash Bargaining Solution. Motivated by recent evidence that panels of LLM
judges lose most of their nominal independence to correlated errors
\citep{kohli2026}, we add a \emph{correlated-error correction}: the arbitration
stage discounts a redundant panel of $n$ evaluators to an effective count
$\neff = \mathbf{1}^\top \Sigma^{-1}\mathbf{1}$ and aggregates their beliefs in
log-odds space, so a correlated panel is not treated as $n$ independent votes.
On seeded synthetic panels with a known error-correlation $\rho$, this
correction is exactly the identity at $\rho=0$ and, at $\rho=0.9$ with $n=3$,
reduces expected calibration error from $0.391$ to $0.267$ and overconfidence
on wrong verdicts from $0.314$ to $0.169$ (Section~\ref{sec:correlation}).
We also contribute a full open-source implementation, a reproducible
correctness test suite, and an ablation harness. \textbf{The end-to-end
LLM-based benchmark (accuracy against human preferences, position-flip rate,
verbosity bias) has not yet been run on live models; those numbers are reported
as \todo{pending real run} rather than estimated,} and the released code
contains no fabricated results. We describe the experimental protocol
(a stratified pilot and a full multi-backbone run on public preference data)
needed to fill them in.
\end{abstract}

\section{Introduction}
Evaluating the output quality of large language models (LLMs) has become a
primary bottleneck in AI development. While human evaluation remains the gold
standard, its cost and lack of scalability have driven the widespread adoption
of LLM-as-a-judge frameworks \citep{zheng2023}. However, single-agent LLM
judges suffer from systematic biases:
\begin{itemize}[leftmargin=1.4em]
  \item \textbf{Position bias:} swapping the presentation order of two candidate
  answers changes the judge's preference in a substantial fraction of cases,
  typically favouring whichever answer appears first \citep{zheng2023}.
  \item \textbf{Verbosity bias:} judges tend to assign higher scores to longer
  answers, even when the extra text is redundant.
  \item \textbf{Sycophancy:} judges tend to agree with a proposer's implicit
  position rather than raising objective criticisms.
\end{itemize}

Prior methods alleviate these biases with prompt engineering, temperature
tuning, or heuristic voter ensembles, but these address symptoms rather than
the root cause: single-agent judges lack the \emph{incentives} to be objective,
concise, and calibrated. A parallel line of work moves to \emph{panels} of
judges \citep{verga2024}, but recent evidence shows that panels sharing base
models make the same mistakes on the same items, so nine judges can carry only
about two independent votes' worth of information \citep{kohli2026}. This
motivates treating evaluator agreement as strategically- and
statistically-structured rather than as independent evidence.

In this work we introduce GT-Judge, a game-theoretic multi-agent evaluation
framework that replaces subjective, single-shot judging with a structured
pipeline whose verdict is the outcome of three coupled games:
\begin{enumerate}[leftmargin=1.6em]
  \item \textbf{Reasoning Debate Game:} a perfect-information extensive-form
  game between a Proposer and a Challenger; a token penalty $\lambda$ is
  deducted from payoffs to discourage verbosity.
  \item \textbf{Confidence Calibration Game:} evaluators independently report a
  probability $\hat{p}\in[0,1]$; payoffs use the Brier reward, a strictly proper
  scoring rule that penalises overconfidence.
  \item \textbf{Consensus Arbitration Game:} disagreements are resolved through
  the Nash Bargaining Solution, weighting each evaluator by debate performance
  and calibrated confidence, and---new in this version---correcting for
  correlated evaluator errors.
\end{enumerate}

\paragraph{Contributions.}
\begin{itemize}[leftmargin=1.4em]
  \item \textbf{Game-theoretic orchestration.} We formulate LLM evaluation as a
  three-stage strategic game coupling extensive-form debate, Brier calibration,
  and Nash bargaining.
  \item \textbf{Theoretical guarantees.} We prove that concise reporting is the
  unique Subgame-Perfect Nash Equilibrium under strict tie-breaking
  (Theorem~\ref{thm:spne}), and prove position invariance
  (Proposition~\ref{prop:position}).
  \item \textbf{Correlated-error correction.} We formalise an effective-evaluator
  count $\neff$ and a log-odds aggregation rule that discounts redundant panels
  (Definition~\ref{def:neff}, Proposition~\ref{prop:correction}), directly
  addressing the independence deficit reported by \citet{kohli2026}, and we
  validate it on seeded synthetic panels (Section~\ref{sec:correlation}).
  \item \textbf{Open, reproducible implementation.} We release the full
  pipeline, a 54-test correctness suite, an ablation harness, and a cost-gated
  experiment runner. We report unmeasured quantities as \todo{pending real run}
  rather than estimating them.
\end{itemize}

\section{Related Work}
\paragraph{LLM-as-a-judge.} Direct prompting of LLMs as evaluators was
popularised by MT-Bench \citep{zheng2023}, which also catalogued position,
verbosity, and self-enhancement biases. Subsequent work studies the
consistency of such judges: TrustJudge \citep{trustjudge2025} characterises
score--comparison inconsistencies and proposes entropy-preserving probabilistic
scoring. Trust-or-Escalate \citep{jung2025} takes a \emph{selective} view:
it estimates the confidence that a human would agree with a judgement (via
Simulated Annotators) and abstains or escalates to a stronger model to give a
provable human-agreement guarantee. GT-Judge differs in that it produces a
verdict without abstention and derives evaluator weights and confidence from an
explicit multi-agent game rather than from single-model confidence.

\paragraph{Judge panels and correlated errors.} PoLL \citep{verga2024} shows
that panels of diverse, smaller models can beat a single large judge on
average. \citet{kohli2026} qualify this sharply: measuring a nine-judge panel
with the Kish effective sample size, they find only $\approx 2$ effective
independent votes, an $8$--$22$ point accuracy shortfall relative to the
independent-voting ideal, and that neither more judges nor smarter aggregation
closes more than a fraction of the gap---the bottleneck is correlated judges,
not the aggregation algorithm. We take this as a design constraint rather than
a target to beat: our correction (Section~\ref{sec:correlation}) does not claim
to recover the lost independent \emph{information}; it uses $\neff$ to keep the
panel's \emph{confidence} honest, so a redundant panel does not behave as though
it had $n$ independent votes.

\paragraph{Game theory in LLM multi-agent systems.} ECON \citep{xie2025} models
cooperative multi-LLM \emph{reasoning} as an incomplete-information game and
uses reinforcement learning to coordinate agents toward a Bayesian Nash
Equilibrium, improving accuracy on math and planning tasks. GT-Judge instead
targets adversarial \emph{evaluation} (judging), is entirely training-free, and
uses Nash \emph{bargaining} over verdicts rather than equilibrium coordination
over answers.

\paragraph{Peer prediction and proper scoring.} Proper scoring rules such as
the Brier score, and peer-prediction mechanisms more broadly, elicit truthful
beliefs without direct access to ground truth \citep{prelec2004}. We adapt
these ideas to online evaluator calibration via a consensus-proxy mechanism
that preserves truthfulness incentives under a non-dictatorship condition
(Proposition~\ref{prop:incentive}).

\section{Theoretical Framework}
Let $\pi = (Q, A_1, A_2)$ be an evaluation instance with a question $Q$ and two
candidate answers. GT-Judge runs two symmetric debates (one per answer),
elicits calibrated confidences, and resolves the verdict via bargaining.

\subsection{Game 1: Reasoning Debate Game (Extensive-Form)}
The debate is a perfect-information extensive-form game
$G_1 = \langle N, \mathcal{A}, U, \lambda\rangle$ with players $N=\{P,C\}$
(Proposer, Challenger) and action space
$\mathcal{A}(h)=\{\text{claim, defend, clarify, attack, concede}\}$.

\begin{assumption}[Action classification]
There exists a deterministic format-mapping $\phi:\mathcal{T}\to\mathcal{A}(h)$
mapping tagged text (e.g.\ \texttt{[ACTION: attack]}) to discrete actions. On a
format violation an agent defaults to the weakest action (clarify for $P$,
concede for $C$) and receives a penalty $\epsilon>0$.
\end{assumption}

\begin{definition}[Endogenous claim validity]\label{def:validity}
The validity $v(c)$ of a Proposer claim $c$ is
\[
v(c) = \begin{cases}
1 & \text{if no attack targets } c\\
1 - \mathrm{sev}(a^*,c) & \text{if } a^* \text{ targets } c \text{ and } P \text{ does not defend}\\
v_0(c) & \text{if } a^* \text{ targets } c \text{ and } P \text{ successfully defends}
\end{cases}
\]
where $\mathrm{sev}(a^*,c)\in[0,1]$ is the attack severity assessed by an
impartial validity assessor and $v_0(c)\in(0,1]$ is the reduced validity after
defense.
\end{definition}

\paragraph{Length-normalised debate score (reconciled with the implementation).}
Let $C_X$ be the claims for answer $X$ and $A_X$ the attacks against them. To
avoid re-introducing the very verbosity bias the game is meant to remove, the
released implementation scores a debate by the \emph{mean} surviving-claim
validity net of attack severity, rather than an unnormalised sum:
\begin{equation}\label{eq:debatescore}
\mathrm{DebateScore}(X) =
\frac{1}{|C_X|}\!\sum_{c\in C_X}\! v(c)
\;-\; \frac{1}{|A_X|}\!\sum_{a\in A_X}\!\mathrm{sev}(a,c(a))
\;-\; \lambda\, \mathrm{Tokens}(X),
\end{equation}
with the convention that an empty sum is $0$. The per-agent utilities used for
equilibrium analysis are the corresponding proposer/challenger terms of
Eq.~\eqref{eq:debatescore}. (An earlier draft wrote this score as an
unnormalised sum; we correct it here to match the code, which normalises so
that a longer debate cannot inflate a score merely by making more claims.)

\begin{assumption}[Strict tie-breaking]
At indifference (an active move's utility equals its expected token cost),
players choose the passive action (Proposer abandons/clarifies, Challenger
concedes).
\end{assumption}

\begin{theorem}[Uniqueness of SPNE]\label{thm:spne}
Under Assumptions 1--2 and finite rounds $K$, backward induction yields a unique
Subgame-Perfect Nash Equilibrium $\sigma^*=(\sigma_P^*,\sigma_C^*)$ with
\[
\sigma_P^*(h)=\begin{cases}\text{defend } c & v(c) > \lambda\,\mathbb{E}[|\text{tokens(defense)}|]\\ \text{abandon } c & \text{otherwise}\end{cases}
\quad
\sigma_C^*(h)=\begin{cases}\text{attack } c & \mathrm{sev}(a^*,c) > \lambda\,\mathbb{E}[|\text{tokens(attack)}|]\\ \text{concede} & \text{otherwise}\end{cases}
\]
\end{theorem}
\begin{proof}
The game is finite-depth and of perfect information. Under Assumption~2 the
payoff difference between any active and passive action is strictly non-zero
except at the boundary, which the weak inequality resolves to the passive
choice. By Zermelo's theorem, backward induction on a finite perfect-information
game with a unique optimal action at every decision node yields a unique SPNE.
\end{proof}

\subsection{Game 2: Confidence Calibration Game}
Each evaluator $i\in\{1,\dots,n\}$ reports a confidence $\hat p_i\in[0,1]$ that
answer A is superior, scored by the Brier reward
\[
U_{\mathrm{cal}}(\hat p_i, y) = 1 - (y-\hat p_i)^2,\qquad y\in\{0,0.5,1\}.
\]
Because ground-truth labels are unavailable at deployment, we use a consensus
proxy.
\begin{definition}[Consensus proxy]\label{def:proxy}
In deployment the ground truth $y$ is replaced by $\tilde y$ derived from the
final verdict $v^*$: $\tilde y=1$ if ``A wins'', $0$ if ``B wins'', $0.5$ if
``Tie''.
\end{definition}
\begin{proposition}[Incentive preservation]\label{prop:incentive}
Under $\tilde y$, reporting the true posterior $\hat p_i^*=\mathbb{E}[v^*\mid s_i]$
is strictly utility-maximising for each agent $i$, provided no single agent can
unilaterally dictate $v^*$ (non-dictatorship).
\end{proposition}
\begin{proof}
The Brier reward is strictly proper, so agent $i$ maximises
$\mathbb{E}[1-(\tilde y-\hat p_i)^2]$ iff $\hat p_i=\mathbb{E}[\tilde y\mid s_i]$.
Since $v^*$ is the joint weighted bargaining outcome of all $n\ge 2$ agents,
non-dictatorship holds and $i$ cannot move $\tilde y$ to match its own report.
\end{proof}

\subsection{Game 3: Consensus Arbitration Game}
Disagreements are resolved by the Nash Bargaining Solution over
$v\in\{\text{A wins, B wins, Tie}\}$:
\begin{equation}\label{eq:nash}
v^* = \arg\max_v \prod_{i=1}^{n}\big(U_i(v)-d_i\big)^{w_i}.
\end{equation}
Here $d_i=-C_H-\beta\cdot\mathrm{Effort}_i$ is the disagreement point (the
human-review fallback cost) and $w_i$ is a bargaining weight from debate and
calibration performance. The per-verdict utility used in the implementation is
\begin{equation}\label{eq:utility}
U_i(v) = \underbrace{1-|y_v-\hat p_i|}_{\text{belief alignment}}
\;+\; \tfrac{1}{2}\,U_{\mathrm{cal},i}
\;+\; \alpha_{\mathrm{deb}}\,\big(1-|y_v-\mathrm{Prior}(v)|\big),
\end{equation}
where $y_v\in\{0,0.5,1\}$ encodes the verdict and $\mathrm{Prior}(v)$ is a
debate-derived prior. The two extra terms beyond raw belief alignment---a
calibration-quality bonus $\tfrac12 U_{\mathrm{cal},i}$ and a debate-prior
alignment term weighted by $\alpha_{\mathrm{deb}}$---were present in the code but
omitted from an earlier draft; we state them explicitly here so the paper and
implementation agree.

\begin{proposition}[Position invariance]\label{prop:position}
Let $\pi=(Q,A,B)$ and $\pi'=(Q,B,A)$. If debates are symmetric and scores are
aggregated across role-swapped rounds, then $\text{GT-Judge}(\pi)=\text{GT-Judge}(\pi')$.
\end{proposition}
\begin{proof}
Both answers undergo symmetric debates, so $\mathrm{DebateScore}(A\mid\pi)=\mathrm{DebateScore}(A\mid\pi')$.
Calibration input is order-independent and evaluator order is counterbalanced,
so posteriors swap symmetrically; the Nash Bargaining Solution is symmetric
under coordinate swaps, hence $v^*(\pi)=v^*(\pi')$.
\end{proof}

\section{Correlated-Error Correction}\label{sec:correlation}
\citet{kohli2026} show that a panel's nominal size $n$ badly overstates its
information content when judges share errors. We take this as a constraint and
correct the arbitration stage's \emph{confidence}, not its accuracy.

\begin{definition}[Effective evaluator count]\label{def:neff}
Let $\Sigma\in\mathbb{R}^{n\times n}$ be the correlation matrix of the
evaluators' errors (estimated from their reported confidences across items, or,
for a single item, assumed equicorrelated with parameter $\rho$). The effective
number of independent evaluators is
\begin{equation}
\neff = \mathbf{1}^\top \Sigma^{-1}\mathbf{1},
\end{equation}
clipped to $[1,n]$. For the equicorrelation matrix
$\Sigma_{ij}=\rho$ $(i\neq j)$, $\Sigma_{ii}=1$, this has the closed form
$\neff = n/\big(1+(n-1)\rho\big)$, so $\neff=n$ at $\rho=0$ and $\neff\to 1$ as
$\rho\to 1$.
\end{definition}

\begin{proposition}[Redundancy-aware pooling]\label{prop:correction}
Let $\ell_i=\mathrm{logit}(\hat p_i)$ and let $\bar\ell$ be their
$\Sigma$-optimal (GLS) weighted mean, and let $\ell_0=\mathrm{logit}(\mathrm{prior})$
be the debate prior in log-odds. Define the corrected aggregate
\begin{equation}\label{eq:pool}
\ell^\star = \ell_0 + \neff\,(\bar\ell-\ell_0),\qquad
\hat p^\star=\sigma(\ell^\star).
\end{equation}
Then (i) at $\rho=0$, $\neff=n$ and Eq.~\eqref{eq:pool} equals the standard
independence-assuming logarithmic opinion pool; (ii) as $\rho\to 1$, $\neff\to1$
and the panel is counted as a single judge. Thus the correction interpolates
between full independent pooling and one effective judge as redundancy grows.
\end{proposition}
\begin{proof}[Proof sketch]
A logarithmic opinion pool of $n$ conditionally-independent judges adds $n$
log-odds increments to the prior: $\ell_0+n(\bar\ell-\ell_0)$. Replacing the
count $n$ by $\neff$ rescales the total evidence to $\neff$ independent
increments. Substituting the equicorrelation closed form of
Definition~\ref{def:neff} gives $\neff=n$ at $\rho=0$ (recovering the pool) and
$\neff\to1$ at $\rho\to1$ (a single judge).
\end{proof}

\paragraph{Synthetic validation.} We validate Eq.~\eqref{eq:pool} on seeded
synthetic panels ($n=3$, 20{,}000 trials, seed 20260718) whose errors are
correlated with a \emph{known} $\rho$; this uses no LLM calls, so the numbers are
reproducible now (\texttt{scripts/validate\_correlation\_correction.py}). The
implemented $\neff$ matches the closed form to $<10^{-6}$ for all tested
$(n,\rho)$. Table~\ref{tab:corr} shows that the correction is exactly the
identity at $\rho=0$ and improves calibration as $\rho$ grows: at $\rho=0.9$ it
cuts expected calibration error (ECE) by $32\%$ and overconfidence on wrong
verdicts by $46\%$. Because the naive log-odds pool is genuinely overconfident
under correlation, Brier also improves; the correction does \emph{not} claim to
recover lost accuracy, consistent with \citet{kohli2026}.

\begin{table}[t]
\centering
\caption{Correlated-error correction on synthetic panels ($n=3$, seeded, no
LLM). Lower is better for all three metrics. ``ocw'' is mean overconfidence on
wrong verdicts. At $\rho=0$ the correction is the identity map.}
\label{tab:corr}
\begin{tabular}{cccccccc}
\toprule
$\rho$ & $\neff$ & ECE (naive) & ECE (corr.) & ocw (naive) & ocw (corr.) & Brier (naive) & Brier (corr.)\\
\midrule
0.0 & 3.00 & 0.389 & 0.389 & 0.213 & 0.213 & 0.081 & 0.081\\
0.3 & 1.88 & 0.380 & 0.330 & 0.260 & 0.193 & 0.124 & 0.120\\
0.6 & 1.36 & 0.392 & 0.298 & 0.291 & 0.179 & 0.153 & 0.142\\
0.9 & 1.07 & 0.391 & 0.267 & 0.314 & 0.169 & 0.177 & 0.158\\
\bottomrule
\end{tabular}
\end{table}

\section{System Implementation and Verification}
GT-Judge is a modular Python pipeline supporting OpenAI, Anthropic, and Google
backbones behind one client interface, with exponential-backoff retries and
temperature-0 response caching. Debates for the two answers run concurrently;
per-run identifier counters (rather than shared instance state) make the
concurrent path free of race conditions. The step-by-step procedure is given in
Algorithm~\ref{alg:pipeline}.

\begin{algorithm}[t]
\caption{GT-Judge Evaluation Pipeline}
\label{alg:pipeline}
\begin{algorithmic}[1]
\Require Question $Q$, answers $A,B$, rounds $K$, penalty $\lambda$
\Ensure Verdict $v^*\in\{\text{A wins, B wins, Tie}\}$
\State Run symmetric debate for $A$ and for $B$ (role-swapped), concurrently
\State Compute validity $v(c)$, severity $\mathrm{sev}(a,c)$, and
$\mathrm{DebateScore}$ via Eq.~\eqref{eq:debatescore}
\For{each evaluator $i$}
  \State Counterbalance presentation order; elicit $\hat p_i\in[0,1]$
\EndFor
\State Estimate $\Sigma$, $\neff$; form corrected aggregate via Eq.~\eqref{eq:pool}
\Repeat
  \State Solve Nash bargaining Eq.~\eqref{eq:nash}; resolve proxy $\tilde y$;
  update $w_i$
\Until{$v^*$ converges or max iterations}
\State \Return $v^*$
\end{algorithmic}
\end{algorithm}

\paragraph{Correctness suite and ablations.} The release includes 54 unit tests
covering the game logic, the multi-provider client, robust JSON parsing, the
$\neff$ identity and GLS weights, and edge cases (single evaluator, ties, empty
answers, malformed model output). An ablation harness runs the full pipeline and
three lesioned variants (no debate, no calibration, no arbitration). We note one
qualitative check that already holds offline: with arbitration removed, the
pipeline loses position invariance (flip rate jumps to $100\%$ on the scripted
harness), whereas the full pipeline does not---consistent with
Proposition~\ref{prop:position}.

\section{Experimental Protocol and Status}
\paragraph{What is \emph{not} yet measured.} An earlier version of this paper
reported $100\%$ accuracy, a $0.0\%$ position-flip rate, and a $0.984$ mean
Brier reward on a five-item set. Those figures were produced by a scripted mock
that was hard-coded to select the answer the dataset labelled correct; they
measure nothing about model behaviour and have been removed. On a balanced,
position-randomised version of the pilot set, the same scripted mock scores only
$20\%$, which is exactly the point: the old numbers were an artefact of the data
and the mock, not evidence. We therefore report the end-to-end evaluation
metrics as \todo{pending real run} and describe the protocol to obtain them.

\paragraph{Tier 1 --- pilot.} A stratified $20$-item pilot (four items each in
coding, math \& logic, writing \& reasoning, summarization, argumentation) with
the superior answer placed in position A or B by a fixed seed and a balanced
$8$/$8$/$4$ split of A/B/Tie labels, so that neither ``always A'' nor ``always
B'' beats $40\%$. Single run on a free-tier backbone. Status: dataset released;
run \todo{pending backbone access}.

\paragraph{Tier 2 --- full.} $150$--$300$ items drawn from public preference
data (e.g.\ MT-Bench / Chatbot Arena \citep{zheng2023}), at least two judge
backbones, $3$--$5$ repeats reported as mean$\pm$std, with the ablations above.
The runner (\texttt{run\_benchmark.py}) estimates and confirms API spend before
issuing any paid call. Status: \todo{pending Tier-2 budget}. The only live
constraint we hit while preparing this release was that the sandbox used for
development has no network route to any external LLM provider, so live runs must
be executed in an environment with provider access.

\section{Discussion}
\paragraph{Computational cost.} GT-Judge requires roughly $2K$ debate calls plus
$n$ calibration calls per item versus one call for a single judge; early stopping
via \emph{concede}, temperature-0 caching, and parallelism mitigate this.
\paragraph{Correlated priors.} Persona diversity (Correctness, Clarity, Depth)
reduces---but does not eliminate---shared error, which is precisely why the
$\neff$ correction of Section~\ref{sec:correlation} operates on top of it rather
than instead of it.

\section{Limitations}
First, the end-to-end LLM benchmark has not been run; the headline accuracy,
flip-rate, and verbosity numbers are pending and are not estimated. Second, the
$\neff$ correction is validated on synthetic panels with a known $\rho$; its
benefit on live judge panels, and the best way to estimate $\Sigma$ online,
remain to be measured. Third, when no per-item correlation estimate is available
the method falls back to an equicorrelation prior $\rho$, a hyperparameter.
Fourth, the multi-provider client has been exercised with an OpenAI-style
interface and unit tests; the Anthropic and Google paths are lazily imported and
untested against live endpoints here. Finally, the formulation is pairwise;
listwise ranking and continuous grading are future work.

\section{Conclusion}
We presented GT-Judge, a training-free game-theoretic multi-agent framework that
targets position bias, verbosity bias, and sycophancy by structuring evaluation
as coupled debate, calibration, and arbitration games, and that corrects for
correlated evaluator errors via an effective-evaluator count. The theory,
implementation, correctness suite, ablation harness, and a reproducible
synthetic validation of the correlation correction are complete; the end-to-end
LLM benchmark is specified but not yet run, and is reported as such.

\section*{Reproducibility and Use of AI Assistance}
All code, configuration, the pilot dataset, the test suite, and the scripts that
produce Table~\ref{tab:corr} are released together; the synthetic validation is
seeded (seed 20260718) and reproducible without any API access. No experimental
number in this paper was generated by a language model or estimated: quantities
that require live model calls are marked \todo{pending real run}. AI coding
assistants were used to help write and refactor code and to draft prose; the
author reviewed and is responsible for all content. (The precise
venue-appropriate wording of AI-use disclosure should be re-checked against the
chosen venue's policy before submission.)

\begin{thebibliography}{9}
\bibitem[Jung et al.(2025)]{jung2025} Jaehun Jung, Faeze Brahman, and Yejin Choi.
Trust or Escalate: LLM Judges with Provable Guarantees for Human Agreement.
\emph{International Conference on Learning Representations (ICLR)}, 2025.
arXiv:2407.18370.

\bibitem[Kohli(2026)]{kohli2026} Guneet Kohli. Nine Judges, Two Effective Votes:
Correlated Errors Undermine LLM Evaluation Panels. Apple Machine Learning
Research, 2026. arXiv:2605.29800.

\bibitem[Prelec(2004)]{prelec2004} Dra\v{z}en Prelec. A Bayesian Truth Serum for
Subjective Data. \emph{Science}, 306(5695):462--466, 2004.

\bibitem[TrustJudge(2025)]{trustjudge2025} TrustJudge: Inconsistencies of
LLM-as-a-Judge and How to Alleviate Them. Preprint, 2025. arXiv:2509.21117.

\bibitem[Verga et al.(2024)]{verga2024} Pat Verga et al. Replacing Judges with
Juries: Evaluating LLM Generations with a Panel of Diverse Models. Preprint,
2024. arXiv:2404.18796.

\bibitem[Xie et al.(2025)]{xie2025} Yi Xie, Zhanke Zhou, Chentao Cao, Qiyu Niu,
Liang Liu, and Bo Han. From Debate to Equilibrium: Belief-Driven Multi-Agent LLM
Reasoning via Bayesian Nash Equilibrium. \emph{International Conference on
Machine Learning (ICML)}, 2025. arXiv:2506.08292.

\bibitem[Zheng et al.(2023)]{zheng2023} Lianmin Zheng, Wei-Lin Chiang, Ying
Sheng, Siyuan Zhuang, Zhanghao Wu, Yonghao Zhuang, Zi Lin, Zhuohan Li, Dacheng
Li, Eric P. Xing, Hao Zhang, Joseph E. Gonzalez, and Ion Stoica. Judging
LLM-as-a-Judge with MT-Bench and Chatbot Arena. \emph{Advances in Neural
Information Processing Systems (NeurIPS)}, Datasets and Benchmarks Track, 2023.
arXiv:2306.05685.
\end{thebibliography}

\end{document}
